\begin{abcliste}
\abc Mit 1~Zoll $= \zZoO$ (FoTaBe S.~185, inch) ist der Radius der Platte
\begin{align}
 r &= \rFpF = \frac{\nZo\cdot\zZo}{2} = \rFp
\end{align}
Die Frequenz ist $f = \nFpO = \fFpP{0}{3}$. Ein Punkt am Rand legt pro Umdrehung den Umfang $2\pi r$ zurück:
\begin{align}
 v &= \vFpF = 2\pi\cdot\rFp\cdot\fFp = \vFp \approx \result{\vFpP{0}{2}}
\end{align}

\abc Im ungünstigsten Fall ist der gesuchte Punkt gerade am Lesekopf vorbeigelaufen und muss eine ganze Umdrehung abwarten. Die maximale Wartezeit ist die Umlaufzeit
\begin{align}
 T &= \TFpF = \frac{1}{\fFp} = \TFp \approx \result{\TFpP{-3}{2}}
\end{align}
Sie hängt nicht vom Abstand zur Plattenmitte ab: Alle Punkte der Platte drehen sich mit derselben Winkelgeschwindigkeit und haben dieselbe Umlaufzeit. Aussen ist nur die Bahngeschwindigkeit grösser.
\end{abcliste}
